\section{Evidence of an emerging mechanics of learning}
\label{sec:reasons_why}


A great cause for optimism that a mechanics of learning is possible is the fact that the essential ingredients of deep learning are both \textit{explicit} and \textit{measurable.}
A deep learning system is characterized by the following components:
\begin{equation*}
    \begin{aligned}
    \text{Architecture:} \quad & \text{a neural network } f(\vx;\vtheta) \text{ specified as a composition of simple linear and nonlinear transformations.} \\
    \text{Data:} \quad & \text{a dataset } \mathcal{D} = \{(\vx_i, \vy_i)\}_{i=1}^n \text{ consisting of samples from an unknown data-generating distribution }\\[-0.7ex]& (\vx,\vy) \sim \mathcal{P}_{\mathrm{data}}. \\
    \text{Task:} \quad & \text{an objective } \mathcal{L}(\vtheta) \text{ measuring the performance of the network } f(\vx;\vtheta) \text{ on the dataset } \mathcal{D}. \\
    \text{Learning rule:} \quad & \text{a gradient-based update equation, e.g. }
    \vtheta^{(t+1)} = \vtheta^{(t)} - \eta \nabla \mathcal{L}(\vtheta^{(t)})
    \text{, together with a parameter} \\[-0.7ex]
    & \text{initialization, e.g. } \vtheta^{(0)}_i \sim \mathcal{N}(0,\alpha^2_{\mathrm{init}}) \text{ and optimization hyperparameters, e.g. learning rate }\eta\text{.}
    \end{aligned}
\end{equation*}
Nothing about the learning process is hidden.
Unlike many complex systems where the equations governing dynamics must be inferred from observations, deep learning directly exposes its ``equations of motion.''
Moreover, these dynamics are extraordinarily measurable: every weight, activation, gradient, and loss value can be recorded, along with arbitrary statistics derived from them.
As a result, deep learning experiments are unusually easy to design, replicate, and interrogate, making it more straightforward to discover empirical regularities and rigorously test theoretical predictions.
Few fast-moving scientific domains offer comparable transparency in their governing equations or comparable freedom in what can be measured.


What, then, stands in the way of a scientific theory of deep learning? The central challenge is not \emph{opacity}, but \emph{complexity}.
While we have direct access to the architecture, data, task, and learning rule, the interaction of these components gives rise to learning dynamics that are nonlinear, coupled, and high-dimensional.
These dynamics depend in subtle ways on the choice of hyperparameters.
And even though we can inspect every training sample, data distributions are complex and have defied simple characterization.


Nevertheless, we argue that this complexity conceals underlying regularities, and that deep learning will indeed admit a scientific theory.
In what follows, we present five broad observations that serve as evidence for an emerging mechanics of learning.
Each of these admits direct analogies to tools and ideas in other disciplines of mechanics.
These are summarized in \cref{tab:lm_physics_comparison}.


\newcommand{\celllist}[1]{%
  \parbox[t]{\linewidth}{\vspace{-2.3mm}#1}%
}

\begin{table}[H]
\centering
\renewcommand{\arraystretch}{1.25}
\setlength{\tabcolsep}{6pt}

\begin{tabular}{
    >{\centering\arraybackslash}p{0.07\textwidth}
    >{\RaggedRight\arraybackslash}p{0.21\textwidth}
    >{\RaggedRight\arraybackslash}p{0.29\textwidth}
    >{\RaggedRight\arraybackslash}p{0.31\textwidth}
}
\toprule
\textbf{Section} & \textbf{Approach} & \textbf{Examples in deep learning} & \textbf{Examples from physics%
} \\
\midrule
\ref{sec:reason-dynamics}
&
solvable settings
&
\celllist{
deep linear networks,\\
kernel regression,\\
multi-index models
}
&
\celllist{
harmonic oscillator,\\
hydrogen atom,\\
Ising model
}
\\[3em]

\ref{sec:reason-limits}
&
simplifying limits
&
\celllist{
lazy vs.\ rich learning,\\
width, depth $\rightarrow \infty$,\\
small initialization
}
&
\celllist{
thermodynamic limit $(n, V \rightarrow \infty)$,\\
classical limit $(\hbar \rightarrow 0)$,\\
hydrodynamic limit $(\vk, \omega \rightarrow 0)$
}
\\[3em]

\ref{sec:reason-measurable}
&
simple empirical laws
&
\celllist{
neural scaling laws\\
edge of stability\\
neural feature ansatz
}
&
\celllist{
the laws of Kepler, Snell, Boyle, Hooke, Newton, Faraday, Ohm, Poiseuille, Planck, Hubble, etc.
}
\\[3em]

\ref{sec:reason-hps}
&
\celllist{
study of system\\
parameters
}
&
\celllist{
step size as sharpness\\
regularization,\\
$\mu$P and width-scaling
}
&
\celllist{
scaling analysis,\\
nondimensionalization,\\
chaotic vs. ordered regimes
}
\\[3em]

\ref{sec:reason-universal-behavior}
&
universal phenomena
&
\celllist{
common inductive biases and representations across models
}
&
\celllist{
critical phenomena,\\
renormalization group flow
}
\\[2em]

\bottomrule
\end{tabular}

\caption{
\textbf{Useful tools and ideas in the emerging science of deep learning closely resemble important tools and ideas from physics, particularly classical mechanics, continuum mechanics, statistical mechanics, and quantum mechanics.}
Extrapolating, this suggests that there will be a \textit{mechanics of learning} which offers a unifying first-principles theory of the training process, hidden representations, final weights, and test-time performance of neural networks.
}
\label{tab:lm_physics_comparison}
\end{table}

